%================================================================
\chapter{Introducción a \LaTeX{}}
\label{ch:latex}
%================================================================

\LaTeX{} es un sistema de preparación de documentos en el que el autor
escribe texto plano intercalado con comandos de marcado, y un
compilador transforma ese código fuente en un PDF compuesto
tipográficamente. Esto contrasta con procesadores de texto como
Microsoft Word o Google Docs, donde el formato se aplica de manera
interactiva en un editor de tipo «lo que ve es lo que obtiene»
(WYSIWYG). La propiedad clave que convierte a \LaTeX{} en el estándar
en ciencia y matemáticas es su capacidad de representar notación
matemática arbitrariamente compleja con precisión y consistencia. Cada
documento producido en este programa es un documento \LaTeX{}.

%================================================================
\section{Su primer documento \LaTeX{}}
\label{sec:first-latex-doc}
%================================================================

Un archivo fuente de \LaTeX{} es un archivo de texto plano con la
extensión \texttt{.tex}. Así como \texttt{hello.py} era un archivo de
texto plano que Python interpretaba, \texttt{hello.tex} es un archivo
de texto plano que el compilador de \LaTeX{} procesa para producir un
PDF. Se aplican las mismas reglas sobre editores de texto (véase la
Sección~\ref{sec:plaintext-first}): utilice Notepad, TextEdit en modo
de texto plano o \texttt{gedit}; nunca un procesador de texto.

\begin{infobox}{Ejercicio L-1: Su primer documento \LaTeX{}}
Abra un editor de texto plano (véase la
Tabla~\ref{tab:text-editors}), escriba exactamente lo siguiente y
guarde el archivo como \texttt{hello.tex}:

\begin{lstlisting}[style=inlineStyle]
\documentclass{article}
\begin{document}

Hello, world. This is my first \LaTeX{} document.

Here is some mathematics: $e^{i\pi} + 1 = 0$.

\end{document}
\end{lstlisting}
\end{infobox}

Cada línea del listado cumple un propósito específico:

\begin{itemize}
    \item \verb|\documentclass{article}| declara la clase de
          documento. Para la mayoría de los informes y artículos,
          \texttt{article} es la elección correcta.
    \item \verb|\begin{document}| y \verb|\end{document}| delimitan el
          cuerpo. Todo lo que esté fuera de este par constituye el
          \emph{preámbulo} (se aborda en la
          Sección~\ref{sec:latex-structure}).
    \item El texto entre los delimitadores es el contenido visible.
    \item \verb|$e^{i\pi} + 1 = 0$| es el modo matemático en línea: el
          texto entre signos de dólar se compone como matemáticas.
\end{itemize}

\begin{table}[htbp]
\centering
\caption{Compilación de un documento \LaTeX{} en cada plataforma}
\label{tab:latex-compile}
\footnotesize
\begin{tabular}{ll}
\toprule
\textbf{Plataforma} & \textbf{Comando} \\
\midrule
Windows (Símbolo del sistema o Git Bash) & \texttt{pdflatex hello.tex} \\
macOS (Terminal) & \texttt{pdflatex hello.tex} \\
Linux & \texttt{pdflatex hello.tex} o \texttt{latexmk -pdf hello.tex} \\
\bottomrule
\end{tabular}
\end{table}

Tenga en cuenta que, con frecuencia, se requieren dos pasadas de
compilación para documentos con referencias cruzadas o un índice
general; \texttt{latexmk} se encarga de esto automáticamente.

\vs
\noindent
La ejecución de \texttt{pdflatex hello.tex} produce varios archivos.
El importante es \texttt{hello.pdf}, que puede abrirse en cualquier
visor de PDF. Los archivos \texttt{.aux}, \texttt{.log} y
\texttt{.out} son auxiliares; pueden eliminarse sin problema.

%================================================================
\section{Un ejemplo más completo}
\label{sec:latex-example2}
%================================================================

El Ejercicio L-1 produjo el documento mínimo viable. Esta sección
presenta un documento más realista que demuestra las características
más necesarias en la escritura académica: una jerarquía estructurada
de secciones, un índice general generado automáticamente, citas en el
texto enlazadas a una bibliografía y la gama completa de modos de
presentación matemática. Cada característica mostrada aquí se utiliza
en los documentos producidos en este programa.

\begin{infobox}{Ejercicio L-2: Un documento estructurado con matemáticas y bibliografía}
Cree dos archivos en la misma carpeta. El primero es el documento
principal, \texttt{report.tex}. El segundo es la base de datos
bibliográfica, \texttt{sample.bib}. Ambos archivos deben estar en el
mismo directorio para que la bibliografía se resuelva.

\vs
\noindent
\textbf{Archivo 1: \texttt{report.tex}}

\begin{lstlisting}[style=inlineStyle]
\documentclass{article}

\usepackage{amsmath}   % entornos matematicos extendidos
\usepackage{booktabs}  % tablas profesionales
\usepackage{hyperref}  % enlaces y referencias cruzadas navegables

\title{A Short Report on Mathematical Notation}
\author{Your Name}
\date{\today}

\begin{document}

\maketitle
\tableofcontents
\newpage

\section{Introduction}

This report demonstrates document structure and mathematical
notation in \LaTeX{}. For background on literate programming,
see Knuth~\cite{knuth1984literate}. The \LaTeX{} system itself
is described in~\cite{lamport1994latex}.

\subsection{Motivation}

Scientific writing requires precise notation. \LaTeX{} provides
that precision through its mathematical typesetting engine.

\subsubsection{A Note on Style}

Use display equations for any expression that warrants visual
emphasis, and inline math for symbols embedded in prose.

\section{Mathematical Display Modes}

\subsection{Inline Math}

Inline math is enclosed in single dollar signs: $f(x) = x^2 + 1$.
It flows within the surrounding text without a line break.

\subsection{Display Math (unnumbered)}

The preferred form for unnumbered display equations is the
\texttt{\textbackslash[...\textbackslash]} environment:

\[
    \int_{-\infty}^{\infty} e^{-x^2}\,dx = \sqrt{\pi}
\]

The legacy \texttt{\$\$...\$\$} form produces the same visual
result but is deprecated in \LaTeX{} and should be avoided in
new documents.

\subsection{Numbered Equations}

The \texttt{equation} environment produces a numbered equation
that can be cross-referenced:

\begin{equation}
    \bar{x} = \frac{1}{n} \sum_{i=1}^{n} x_i
    \label{eq:mean}
\end{equation}

Equation~\ref{eq:mean} defines the sample mean. An approximation
for the standard deviation is given in~\cite{gurland1971simple}.

\subsection{Multi-line Aligned Equations}

The \texttt{align} environment aligns multiple equations at a
designated column, marked with \texttt{\&}:

\begin{align}
    \mu    &= \mathbb{E}[X]            \label{eq:mu}   \\
    \sigma^2 &= \mathbb{E}[(X-\mu)^2] \label{eq:var}
\end{align}

Equations~\ref{eq:mu} and~\ref{eq:var} define the mean and
variance of a random variable $X$.

\section{A Simple Table}

\begin{table}[htbp]
\centering
\caption{Comparison of equation display modes}
\label{tab:modes}
\begin{tabular}{lll}
\toprule
Mode & Syntax & Use case \\
\midrule
Inline     & \texttt{\$...\$}                    & Symbols in prose \\
Unnumbered & \texttt{\textbackslash[...\textbackslash]}  & Emphasis, no reference needed \\
Numbered   & \texttt{equation}                   & Cross-referenceable \\
Aligned    & \texttt{align}                      & Multi-line derivations \\
\bottomrule
\end{tabular}
\end{table}

\bibliographystyle{plain}
\bibliography{sample}

\end{document}
\end{lstlisting}

\vs
\noindent
\textbf{Archivo 2: \texttt{sample.bib}}

\begin{lstlisting}[style=inlineStyle]
@article{knuth1984literate,
  author  = {Knuth, Donald E.},
  title   = {Literate Programming},
  journal = {The Computer Journal},
  year    = {1984},
  volume  = {27},
  number  = {2},
  pages   = {97--111}
}

@book{lamport1994latex,
  author    = {Lamport, Leslie},
  title     = {\LaTeX: A Document Preparation System},
  edition   = {2nd},
  publisher = {Addison-Wesley},
  year      = {1994}
}

@article{gurland1971simple,
  author  = {Gurland, John and Tripathi, Ram C.},
  title   = {A Simple Approximation for Unbiased Estimation
             of the Standard Deviation},
  journal = {The American Statistician},
  year    = {1971},
  volume  = {25},
  number  = {4},
  pages   = {30--32}
}
\end{lstlisting}
\end{infobox}

\begin{warningbox}{La compilación con bibliografía requiere cuatro pasos}
Una única pasada de \texttt{pdflatex} no es suficiente cuando el
documento contiene comandos \verb|\cite{}|. La secuencia completa es:

\begin{verbatim}
pdflatex report.tex
bibtex report
pdflatex report.tex
pdflatex report.tex
\end{verbatim}

La primera pasada crea el archivo \texttt{.aux} que enumera todas las
claves de cita. \texttt{bibtex} lee el archivo \texttt{.aux} y la
base de datos \texttt{.bib} para producir un archivo \texttt{.bbl}
que contiene la bibliografía formateada. La segunda y tercera pasadas
de \texttt{pdflatex} incorporan la bibliografía y resuelven todas las
referencias cruzadas. La ejecución de \texttt{latexmk -pdf report.tex}
realiza todas las pasadas necesarias de manera automática.
\end{warningbox}

\begin{table}[htbp]
\centering
\caption{Modos de presentación de ecuaciones en \LaTeX{}}
\label{tab:equation-modes}
\begin{tabular}{p{2.8cm} p{4.2cm} p{2.5cm} p{4.5cm}}
\toprule
\textbf{Modo} & \textbf{Sintaxis} & \textbf{Estilo de salida} & \textbf{Notas} \\
\midrule
En línea & \verb|$...$| & Dentro del texto, compacto & Se utiliza para símbolos y expresiones breves en la prosa \\
Presentación (heredado) & \verb|$$...$$| & Centrado, tamaño completo & Obsoleto; evítese en documentos nuevos \\
Presentación (preferido) & \verb|\[...\]| & Centrado, tamaño completo & Se utiliza para ecuaciones de presentación sin numerar \\
Numerado & \verb|\begin{equation}| \verb|...\end{equation}| & Centrado con número & Referenciable mediante \verb|\label|/\verb|\ref| \\
Alineado de varias líneas & \verb|\begin{align}| \verb|...\end{align}| & Varias líneas, alineadas en \verb|&| & Cada línea puede llevar su propia \verb|\label| \\
\bottomrule
\end{tabular}
\end{table}

%================================================================
\section{Estructura del documento}
\label{sec:latex-structure}
%================================================================

Todo archivo fuente de \LaTeX{} consta de dos partes: el
\emph{preámbulo} (todo lo que antecede a \verb|\begin{document}|) y el
\emph{cuerpo} (todo lo que se encuentra entre \verb|\begin{document}|
y \verb|\end{document}|). El preámbulo declara la clase de documento,
carga paquetes con \verb|\usepackage{...}| y establece los metadatos.
El cuerpo contiene el contenido visible, organizado mediante comandos
de seccionamiento.

\begin{lstlisting}[style=inlineStyle, caption={Documento mínimo compilable con preámbulo y seccionamiento}]
\documentclass{article}

% --- Preambulo ---
\usepackage{amsmath}    % matematicas extendidas
\usepackage{graphicx}   % figuras
\usepackage{booktabs}   % tablas profesionales
\usepackage{hyperref}   % enlaces navegables

\title{My First Report}
\author{Alice Smith}
\date{\today}

% --- Cuerpo ---
\begin{document}

\maketitle

\section{Introduction}
This is the introduction.

\subsection{Background}
Some background material.

\subsubsection{Details}
More specific details.

\section{Methods}
Description of methods.

\end{document}
\end{lstlisting}

El comando \verb|\usepackage{...}| carga paquetes de extensión. Los
paquetes más relevantes para este programa son \texttt{amsmath}
(matemáticas extendidas), \texttt{graphicx} (figuras),
\texttt{booktabs} (tablas) y \texttt{hyperref} (enlaces navegables).

%================================================================
\section{Matemáticas en \LaTeX{}: una hoja de referencia}
\label{sec:latex-math}
%================================================================

\LaTeX{} distingue dos modos matemáticos. El modo en línea
(\verb|$...$|) compone matemáticas dentro de una línea de texto. El
modo de presentación (\verb|\[...\]| o el entorno \texttt{equation})
centra las matemáticas en su propia línea y, al utilizar
\texttt{equation}, asigna un número de ecuación. Utilice el modo de
presentación para cualquier expresión que merezca su propio énfasis
visual, y el modo en línea para símbolos incrustados en la prosa.

%----------------------------------------------------------------
\subsection{Notación matemática común}
%----------------------------------------------------------------

\begin{longtable}{p{3.2cm} p{4.5cm} p{3.0cm} p{3.5cm}}
\caption{Notación matemática común en \LaTeX{}} \label{tab:latex-math-cheatsheet} \\
\toprule
\textbf{Categoría} & \textbf{Código fuente \LaTeX{}} & \textbf{Renderizado} & \textbf{Notas} \\
\midrule
\endfirsthead
\toprule
\textbf{Categoría} & \textbf{Código fuente \LaTeX{}} & \textbf{Renderizado} & \textbf{Notas} \\
\midrule
\endhead
\midrule
\multicolumn{4}{r}{\emph{Continúa en la página siguiente}} \\
\endfoot
\bottomrule
\endlastfoot

% --- Modos en linea y de presentacion ---
\multicolumn{4}{l}{\textbf{Modos en línea y de presentación}} \\
\midrule
Matemática en línea & \verb|$x + y = z$| & $x + y = z$ & Dentro de una oración \\
Matemática en presentación & \verb|\[ x + y = z \]| & (bloque centrado) & Sin número \\
Ecuación numerada & \verb|\begin{equation}...\end{equation}| & (con número) & Referenciable \\
\midrule

% --- Aritmetica basica ---
\multicolumn{4}{l}{\textbf{Aritmética básica}} \\
\midrule
Suma & \verb|a + b| & $a+b$ & \\
Resta & \verb|a - b| & $a-b$ & \\
Multiplicación (punto) & \verb|a \cdot b| & $a \cdot b$ & Preferido en matemáticas \\
Multiplicación (cruz) & \verb|a \times b| & $a \times b$ & Vectores, producto vectorial \\
División & \verb|\frac{a}{b}| & $\frac{a}{b}$ & Utilice siempre \verb|\frac| en presentación \\
Superíndice & \verb|x^{2}| & $x^{2}$ & Se requieren llaves para >1 carácter \\
Subíndice & \verb|x_{i}| & $x_{i}$ & Se requieren llaves para >1 carácter \\
Ambos & \verb|x_{i}^{2}| & $x_{i}^{2}$ & \\
\midrule

% --- Funciones comunes ---
\multicolumn{4}{l}{\textbf{Funciones comunes}} \\
\midrule
Raíz cuadrada & \verb|\sqrt{x}| & $\sqrt{x}$ & \\
Raíz $n$-ésima & \verb|\sqrt[n]{x}| & $\sqrt[n]{x}$ & \\
Exponencial & \verb|e^{x}| o \verb|\exp(x)| & $e^{x}$ & \verb|\exp| preferido en presentación \\
Logaritmo natural & \verb|\ln x| & $\ln x$ & \\
Logaritmo en base $b$ & \verb|\log_{b} x| & $\log_{b} x$ & \\
Seno & \verb|\sin x| & $\sin x$ & Utilice \verb|\sin|, no \texttt{sin} \\
Coseno & \verb|\cos x| & $\cos x$ & \\
Valor absoluto & \verb+|x|+ o \verb|\lvert x \rvert| & $\lvert x \rvert$ & \verb|\lvert| se escala con el contenido \\
Parte entera inferior & \verb|\lfloor x \rfloor| & $\lfloor x \rfloor$ & \\
Parte entera superior & \verb|\lceil x \rceil| & $\lceil x \rceil$ & \\
\midrule

% --- Sumatoria e integracion ---
\multicolumn{4}{l}{\textbf{Sumatoria e integración}} \\
\midrule
Sumatoria & \verb|\sum_{i=1}^{n} x_i| & $\sum_{i=1}^{n} x_i$ & Límites arriba/abajo en presentación \\
Productoria & \verb|\prod_{i=1}^{n} x_i| & $\prod_{i=1}^{n} x_i$ & \\
Integral & \verb|\int_{a}^{b} f(x)\,dx| & $\int_{a}^{b} f(x)\,dx$ & \verb|\,| añade un espacio fino antes de $dx$ \\
Límite & \verb|\lim_{x \to 0} f(x)| & $\lim_{x \to 0} f(x)$ & \\
\midrule

% --- Letras griegas ---
\multicolumn{4}{l}{\textbf{Letras griegas (comunes en bioestadística y aprendizaje automático)}} \\
\midrule
alfa & \verb|\alpha| & $\alpha$ & Nivel de significancia, tasa de aprendizaje \\
beta & \verb|\beta| & $\beta$ & Coeficientes de regresión \\
gamma & \verb|\gamma|, \verb|\Gamma| & $\gamma$, $\Gamma$ & Distribución gamma \\
delta & \verb|\delta|, \verb|\Delta| & $\delta$, $\Delta$ & Cambio, delta de Dirac \\
épsilon & \verb|\epsilon|, \verb|\varepsilon| & $\epsilon$, $\varepsilon$ & Cantidad pequeña \\
theta & \verb|\theta|, \verb|\Theta| & $\theta$, $\Theta$ & Parámetros, ángulo \\
lambda & \verb|\lambda|, \verb|\Lambda| & $\lambda$, $\Lambda$ & Tasa de Poisson, valor propio \\
mu & \verb|\mu| & $\mu$ & Media \\
nu & \verb|\nu| & $\nu$ & Grados de libertad \\
pi & \verb|\pi|, \verb|\Pi| & $\pi$, $\Pi$ & 3{,}14\ldots, productoria \\
rho & \verb|\rho| & $\rho$ & Correlación \\
sigma & \verb|\sigma|, \verb|\Sigma| & $\sigma$, $\Sigma$ & Desviación estándar, sumatoria \\
tau & \verb|\tau| & $\tau$ & Constante de tiempo \\
phi & \verb|\phi|, \verb|\Phi| & $\phi$, $\Phi$ & CDF normal, ángulo \\
chi & \verb|\chi| & $\chi$ & Prueba de chi-cuadrado \\
omega & \verb|\omega|, \verb|\Omega| & $\omega$, $\Omega$ & Frecuencia, espacio muestral \\
\midrule

% --- Relaciones y logica ---
\multicolumn{4}{l}{\textbf{Relaciones y lógica}} \\
\midrule
Distinto de & \verb|\neq| & $\neq$ & \\
Menor/mayor o igual & \verb|\leq|, \verb|\geq| & $\leq$, $\geq$ & \\
Aproximadamente & \verb|\approx| & $\approx$ & \\
Proporcional a & \verb|\propto| & $\propto$ & \\
Distribuido como & \verb|\sim| & $\sim$ & \\
Para todo & \verb|\forall| & $\forall$ & \\
Existe & \verb|\exists| & $\exists$ & \\
Implica & \verb|\Rightarrow| & $\Rightarrow$ & \\
Si y sólo si & \verb|\Leftrightarrow| & $\Leftrightarrow$ & \\
Pertenece a & \verb|\in| & $\in$ & \\
Subconjunto & \verb|\subset|, \verb|\subseteq| & $\subset$, $\subseteq$ & \\
Unión, intersección & \verb|\cup|, \verb|\cap| & $\cup$, $\cap$ & \\
\midrule

% --- Conjuntos numericos ---
\multicolumn{4}{l}{\textbf{Conjuntos numéricos}} \\
\midrule
Números reales & \verb|\mathbb{R}| & $\mathbb{R}$ & Requiere \texttt{amsfonts} o \texttt{amssymb} \\
Números naturales & \verb|\mathbb{N}| & $\mathbb{N}$ & \\
Enteros & \verb|\mathbb{Z}| & $\mathbb{Z}$ & \\
Racionales & \verb|\mathbb{Q}| & $\mathbb{Q}$ & \\
Complejos & \verb|\mathbb{C}| & $\mathbb{C}$ & \\
\midrule

% --- Vectores y matrices ---
\multicolumn{4}{l}{\textbf{Vectores y matrices}} \\
\midrule
Vector en negrita & \verb|\mathbf{x}| & $\mathbf{x}$ & Preferido en estadística \\
Vector con flecha & \verb|\vec{x}| & $\vec{x}$ & Convención física \\
Transpuesta & \verb|A^{\top}| o \verb|A^T| & $A^{\top}$ & \\
Matriz $2\times 2$ & \verb|\begin{bmatrix} a & b \\ c & d \end{bmatrix}| & $\begin{bmatrix} a & b \\ c & d \end{bmatrix}$ & Requiere \texttt{amsmath} \\
Norma & \verb+\|x\|+ & $\|x\|$ & \\
Producto interno & \verb|\langle x, y \rangle| & $\langle x, y \rangle$ & \\

\end{longtable}

%================================================================
\section{Figuras}
\label{sec:latex-figures}
%================================================================

Las figuras se incluyen utilizando el entorno flotante \texttt{figure}
y el comando \verb|\includegraphics| del paquete \texttt{graphicx}. El
compilador \texttt{pdflatex} acepta los formatos de imagen PDF, PNG y
JPG.

\begin{lstlisting}[style=inlineStyle, caption={Inclusión de una figura en \LaTeX{}}]
\begin{figure}[htbp]
    \centering
    \includegraphics[width=0.7\textwidth]{images/my-plot.png}
    \caption{Description of the figure.}
    \label{fig:my-plot}
\end{figure}
\end{lstlisting}

La opción \texttt{[htbp]} sugiere preferencias de ubicación a
\LaTeX{}: aquí (\emph{here}), arriba (\emph{top}), abajo
(\emph{bottom}) o en una página dedicada. Los comandos
\verb|\caption| y \verb|\label| proporcionan una leyenda numerada y un
objetivo de referencia cruzada, respectivamente.

%================================================================
\section{Tablas}
\label{sec:latex-tables}
%================================================================

Las tablas se construyen con el flotante \texttt{table} y el entorno
\texttt{tabular}. El paquete \texttt{booktabs} proporciona líneas
horizontales limpias. Por convención, las líneas verticales
(\texttt{|}) se omiten en los documentos profesionales; las líneas
horizontales de \texttt{booktabs} son suficientes.

\begin{lstlisting}[style=inlineStyle, caption={Una tabla mínima con booktabs}]
\begin{table}[htbp]
\centering
\caption{Sample measurements}
\label{tab:sample}
\begin{tabular}{lrr}
\toprule
Subject & Height (cm) & Weight (kg) \\
\midrule
Alice   & 165         & 60          \\
Bob     & 178         & 75          \\
Carol   & 172         & 68          \\
\bottomrule
\end{tabular}
\end{table}
\end{lstlisting}

%================================================================
\section{Bibliografía}
\label{sec:latex-bibliography}
%================================================================

\LaTeX{} gestiona las bibliografías mediante un archivo externo
(\texttt{.bib}) y un comando de estilo. El comando \verb|\cite{key}|
inserta una cita en cualquier punto del texto; los comandos
\verb|\bibliography{filename}| y \verb|\bibliographystyle{plain}|
generan la lista de referencias.

Para una gestión bibliográfica más sofisticada, el paquete
\texttt{biblatex} con el motor \texttt{biber} es preferido en los
documentos modernos. El preámbulo compartido de este repositorio
tiene el soporte bibliográfico comentado; los estudiantes que
escriban informes independientes deben descomentarlo y configurarlo.

%================================================================
\section{LaTeX Workshop en Visual Studio Code}
\label{sec:latex-workshop}
%================================================================

La extensión LaTeX Workshop para VSC (mencionada en la
Sección~\ref{sec:extensions}) proporciona un flujo de trabajo
integrado completo: resaltado de sintaxis, vista previa de PDF en
vivo, compilación con un solo clic y navegación de errores. Una vez
instalada, al abrir cualquier archivo \texttt{.tex} se activa la
extensión automáticamente. El atajo de compilación es
\texttt{Ctrl+Alt+B} (Windows/Linux) o \texttt{Cmd+Option+B} (macOS),
que ejecuta \texttt{pdflatex} y muestra el resultado en el visor de
PDF integrado. Los estudiantes no necesitan utilizar la línea de
comandos para la compilación después de configurar esto; el enfoque
por línea de comandos de la Sección~\ref{sec:first-latex-doc} se
presenta primero para que se comprenda el proceso subyacente.

%================================================================
\section{Resumen}
\label{sec:latex-summary}
%================================================================

Este capítulo introdujo \LaTeX{}, el sistema de preparación de
documentos utilizado a lo largo de este programa y ampliamente
adoptado como el estándar en ciencia y matemáticas:

\begin{itemize}
    \item \textbf{Qué es \LaTeX{}.} Un lenguaje de marcado de texto
          plano compilado en PDF compuestos tipográficamente, en
          contraste con los procesadores de texto WYSIWYG.
    \item \textbf{El flujo de trabajo desde la fuente hasta el PDF.}
          Escribir un archivo \texttt{.tex} en un editor de texto,
          compilar con \texttt{pdflatex} o \texttt{latexmk} y
          visualizar el PDF resultante.
    \item \textbf{Estructura del documento.} El preámbulo (clase de
          documento, paquetes, metadatos) y el cuerpo (contenido
          organizado con \verb|\section|, \verb|\subsection| y
          \verb|\subsubsection|).
    \item \textbf{Matemáticas en línea y en presentación.} Signos de
          dólar para matemáticas en línea y \verb|\[...\]| o el
          entorno \texttt{equation} para matemáticas en modo
          presentación.
    \item \textbf{Modos de ecuación.} Existen cinco modos: en línea
          (\verb|$...$|), presentación heredada (\verb|$$...$$|,
          obsoleto), presentación sin numerar preferida
          (\verb|\[...\]|), numerado (\texttt{equation}) y alineado de
          varias líneas (\texttt{align}). Utilice \verb|\[...\]| para
          ecuaciones sin numerar y \texttt{equation} para ecuaciones
          de presentación referenciables.
    \item \textbf{La hoja de referencia matemática.} Una tabla de
          referencia exhaustiva (Tabla~\ref{tab:latex-math-cheatsheet})
          que abarca aritmética, funciones, sumatoria, integración,
          letras griegas, relaciones, conjuntos numéricos y
          vectores/matrices.
    \item \textbf{Figuras y tablas.} Entornos flotantes con
          \verb|\includegraphics|, \texttt{tabular} y
          \texttt{booktabs} para un formato profesional.
    \item \textbf{Flujo de trabajo bibliográfico.} Un archivo
          \texttt{.bib} almacena las entradas de referencia. El
          comando \verb|\cite{key}| inserta una cita;
          \verb|\bibliography{filename}| y
          \verb|\bibliographystyle{plain}| generan la lista de
          referencias. Los documentos con citas requieren una
          secuencia de compilación de cuatro pasos: dos pasadas de
          \texttt{pdflatex} que encierran una pasada de
          \texttt{bibtex}, seguida de una \texttt{pdflatex} final.
          \texttt{latexmk} se encarga de esto automáticamente.
    \item \textbf{LaTeX Workshop.} La extensión de VSC que
          proporciona resaltado de sintaxis, vista previa en vivo y
          compilación con un solo clic, eliminando la necesidad de
          compilación por línea de comandos en el uso diario.
\end{itemize}
